% Part: proof-theory
% Chapter: natural-deduction
% Section: introduction

\documentclass[../../../include/open-logic-section]{subfiles}

\begin{document}

\olfileid[id]{pt}{nat}{int}

\olsection{Pendahuluan}

Deduksi alami adalah keluarga sistem !!{proof} yang di dalamnya setiap
penghubung logis atau kuantor memiliki sepasang inferensi: satu yang
``mengintroduksi'' operator tersebut (disebut !!a{introduction}) dan satu
yang ``mengeliminasi'' operator tersebut (disebut !!a{elimination}).
Misalnya, \Intro{\lif} memiliki !!{formula} berbentuk~$!A \lif !B$ sebagai
kesimpulan, sedangkan \Elim{\lif} memiliki !!{formula} berbentuk~$!A \lif
!B$ sebagai premis.

Dua versi pertama deduksi alami diperkenalkan pada tahun 1930-an oleh
Gerhard Gentzen dan Stanisław Jaśkowski. Sistem Gentzen merupakan sistem
yang paling sering dibahas oleh para teoretikus bukti, sedangkan sistem
Jaśkowski melahirkan sistem-sistem yang paling luas digunakan dalam
pengajaran. Dalam kedua sistem itu, kita dapat memulai dengan
\emph{asumsi}---!!{formula} yang tidak perlu di-!!{prove}, tetapi dapat
digunakan dalam !!{proving} !!{formula} lain. Seluruh !!{proof} dapat
bergantung pada asumsi-asumsi semacam itu. Karena asumsi-asumsi tersebut
belum dibuktikan, !!{proof} itu tidak menetapkan kesimpulannya
(\emph{!!{formula} akhir}), melainkan hanya bahwa kesimpulannya mengikuti
dari asumsi-asumsi tersebut.

Sebagian aturan inferensi mempunyai kesimpulan yang tidak lagi bergantung
pada asumsi tertentu: aturan-aturan itu ``melepaskan'' asumsi.
Misalnya, aturan \Intro{\lif} mengambil !!a{proof} atas kesimpulan~$!B$
dari asumsi~$!A$, dan mempunyai $!A \lif !B$ sebagai kesimpulannya.
!!{proof} yang berakhir pada $!A \lif !B$ tidak lagi bergantung
pada~$!A$; asumsi $!A$ telah dilepaskan. Dalam sistem Gentzen,
!!{proof} merupakan pohon !!{formula}; asumsi-asumsinya adalah
!!{formula} paling atas; dan aturan seperti~\Intro{\lif} membawa label
yang menunjukkan asumsi mana yang dilepaskan. Dalam sistem Jaśkowski,
bagian !!{proof} yang bergantung pada suatu asumsi dipisahkan dengan cara
tertentu, misalnya diberi inden di belakang garis vertikal atau ditempatkan
dalam kotak. Asumsi~$!A$ yang dilepaskan dan konsekuen $!B$ dari
kondisional menjadi baris pertama dan terakhir di dalam kotak, sedangkan
kesimpulan~$!A \lif !B$ dari \Intro{\lif} berada di luar kotak.

Sistem-sistem tersebut bersifat ``alami'' karena aturan inferensinya
mencerminkan cara kita (atau setidaknya matematikawan) bernalar secara
alami. Untuk menetapkan hasil hipotetis (suatu kondisional), kita
mengasumsikan antesedennya dan menggunakannya untuk membuktikan
konsekuennya. Itulah~\Intro{\lif}. Ketika kita mengetahui kondisional
$!A \lif !B$ dan antesedennya~$!A$, kita menyimpulkan~$!B$.
Itulah~\Elim{\lif}. Kita menggunakan bukti dengan kasus untuk menunjukkan
bahwa, dengan mengasumsikan disjungsi~$!A \lor !B$, suatu kesimpulan
berlaku. Itulah~\Elim{\lor}. Untuk membuktikan klaim umum
$\lforall[x][!A(x)]$, kita membuktikan bahwa $!A(c)$ berlaku untuk
$c$ yang sebarang. Itulah~\Intro{\lforall}. Demikian seterusnya.

Sebagai contoh, berikut !!{proof} sederhana atas~$!A \lif (!A \lor !B)$
dalam deduksi alami bergaya Gentzen:
\[
\AxiomC{$\Discharge{!A}{x}$}
\RightLabel{\Intro{\lor}}
\UnaryInfC{$!A \lor !B$}
\DischargeRule{\Intro{\lif}}{x}
\UnaryInfC{$!A \lif (!A \lor !B)$}
\DisplayProof
\]
Bukti ini bermula dengan asumsi~$!A$, menggunakan \Intro{\lor} untuk
menyimpulkan~$!A \lor !B$, lalu menggunakan \Intro{\lif} untuk
menyimpulkan~$!A \lif (!A \lor !B)$. Inferensi \Intro{\lif} melepaskan
asumsi~$!A$; hal ini ditunjukkan oleh label~$x$.

Para teoretikus bukti lebih menyukai sistem asli Gentzen yang bekerja
dengan pohon !!{formula}, meskipun salah satu variannya juga penting dalam
teori bukti. !!^a{proof} atas~$!A$ dari asumsi-asumsi~$\Gamma$ menunjukkan
bahwa $!A$ mengikuti dari~$\Gamma$, yakni $\Gamma \Entails !A$. Aturan
inferensi deduksi alami dapat dibaca secara alami sebagai aturan yang
memberi tahu kita sesuatu tentang relasi konsekuensi semacam itu. Misalnya,
\Intro{\lif} mengatakan bahwa jika $\Gamma + !A \Entails !B$, maka
$\Gamma \Entails !A \lif !B$. Wajar pula jika aturan-aturan tersebut
dirumuskan langsung agar bekerja sekaligus pada asumsi dan kesimpulan yang
mengikuti darinya dalam premis serta kesimpulan aturan inferensi; yakni,
aturan itu bekerja pada sekuen $\Gamma \Sequent !A$. Bukti sederhana di
atas akan tampak sebagai berikut dalam sistem semacam itu:
\[
\Axiom$x:!A \fCenter !A$
\RightLabel{\Intro{\lor}}
\UnaryInf$x:!A \fCenter !A \lor !B$
\DischargeRule{\Intro{\lif}}{x}
\UnaryInf$\fCenter !A \lif (!A \lor !B)$
\DisplayProof
\]
Seperti terlihat, aturan-aturannya sama: aturan itu bekerja pada
!!{formula} di sebelah kanan~$\Sequent$, sedangkan !!{formula} di sebelah
kiri hanya mencantumkan asumsi yang menjadi tempat ketergantungan pada
setiap langkah.

Pasangan lengkap aturan !!{introduction}/!!{elimination} itu sendiri
tidak lengkap bagi logika klasik. Kita perlu menambahkan dua aturan yang
berkaitan dengan~$\lfalse$. Yang pertama ialah
``aturan absurditas intuisionistik''~\FalseInt, atau ``ledakan,''
yang mengatakan bahwa dari $\lfalse$ kita dapat menyimpulkan apa pun.
Yang kedua ialah prinsip klasik bukti tak langsung~\FalseCl{} (atau
``aturan absurditas klasik''), yang mengatakan bahwa jika kita dapat
!!{prove}~$\lfalse$ dari~$\lnot !A$, maka kita telah !!{prove}~$!A$.
Jika \FalseCl{} ditiadakan, kita memperoleh sistem yang sesuai untuk
logika intuisionistik. Jika keduanya ditiadakan, kita memperoleh apa yang
dikenal sebagai ``logika minimal.''

Kita akan membahas deduksi alami bergaya Gentzen untuk logika klasik dan
intuisionistik. Sistem-sistem ini adalah \Log{N1c} dan \Log{N1i}.
Sistem \Log{N2c} dan \Log{N2i} merupakan sistem yang bersesuaian, tetapi
bekerja pada sekuen $\Gamma \Sequent !A$, bukan pada satu
!!{formula}~$!A$.

\end{document}
